Mechanistic interpretability has produced substantial progress in understanding the internal representations of large language models, yet its application to Vision Transformers is still emerging. 
Unlike vision-language models such as CLIP, purely visual transformers lack a built-in text-image alignment that can be exploited for grounding discovered features in human-understandable concepts. 
We address this gap by training Sparse Autoencoders (SAEs) on the MLP output activations of a pre-trained ViT-B/16 at multiple network depths, and by using CLIP as an external cross-modal evaluator to assign interpretable labels to the learned sparse directions. 
Across the two monitored layers, the SAEs recover dictionaries of 6,144 features each; for each layer we select and evaluate the $K{=}10$ most active features --- a representative sample rather than an exhaustive census of the full dictionary --- all 10 features per layer receive a distinct CLIP label; visual inspection of the top-$K$ crops confirms monosemantic activation for the causally relevant features.
Mid-network features (Layer~6) are dominated by low-level texture and geometric patterns (\textit{spotted pattern}, \textit{scale pattern}, \textit{honeycomb pattern}), while late-network features (Layer~11) encompass more semantic concepts (\textit{animal}, \textit{plant}, \textit{bird}); causal impact is negligible at Layer~6 (mean logit drop 0.03\%) and concentrated at Layer~11 (mean 1.40\%), driven primarily by a single highly specialised feature (Feature~5065, \textit{honeycomb pattern}, 13.49\% relative logit drop under ablation).
Our results suggest that SAE-based feature decomposition transfers to purely visual transformers and that cross-modal evaluation via an external vision-language model is a viable strategy for bridging the labelling gap that arises in the absence of internal text supervision.
